%%%%%%%%%%%%%%%%%%%%%%%%%%%%%%%%%%%%%%%%%%%%%%%%%%%%%%%%%%%%%%%%%%%%%%%%%%%%%
%% The Minimum Analytic Finite Graph:
%% Truth, Individuation, and the Contact Floor
%% in Finite Weighted Graphs
%%
%% Author: Kundai Farai Sachikonye
%%         Technical University of Munich
%%%%%%%%%%%%%%%%%%%%%%%%%%%%%%%%%%%%%%%%%%%%%%%%%%%%%%%%%%%%%%%%%%%%%%%%%%%%%

\documentclass[12pt, a4paper]{article}

\usepackage{amsmath, amssymb, amsthm}
\usepackage{mathtools}
\usepackage[a4paper, margin=1.15in]{geometry}
\usepackage[colorlinks=true, linkcolor=blue, citecolor=blue, urlcolor=blue]{hyperref}
\usepackage{booktabs}
\usepackage{enumitem}
\usepackage{microtype}
\usepackage{natbib}
\usepackage{graphicx}
\usepackage{float}
\usepackage{tikz}
\usetikzlibrary{graphs, graphdrawing, positioning, arrows.meta}

%% Theorem environments
\theoremstyle{plain}
\newtheorem{theorem}{Theorem}[section]
\newtheorem{lemma}[theorem]{Lemma}
\newtheorem{corollary}[theorem]{Corollary}
\newtheorem{proposition}[theorem]{Proposition}

\theoremstyle{definition}
\newtheorem{definition}[theorem]{Definition}
\newtheorem{example}[theorem]{Example}

\theoremstyle{remark}
\newtheorem{remark}[theorem]{Remark}

%% Notation
\newcommand{\Gfin}{G_{\mathrm{fin}}}
\newcommand{\Kn}{K_n}
\newcommand{\wmin}{w_{\min}}
\newcommand{\cut}{\mathrm{cut}}
\newcommand{\vol}{\mathrm{vol}}
\newcommand{\med}{\mathcal{M}}
\newcommand{\Nbd}{\mathcal{N}}
\newcommand{\reals}{\mathbb{R}}
\newcommand{\naturals}{\mathbb{N}}
\newcommand{\Pcal}{\mathcal{P}}
\newcommand{\Ecal}{\mathcal{E}}
\newcommand{\Vcal}{\mathcal{V}}
\newcommand{\bmin}{\beta_{\min}}
\newcommand{\qed}{\hfill$\blacksquare$}

%%%%%%%%%%%%%%%%%%%%%%%%%%%%%%%%%%%%%%%%%%%%%%%%%%%%%%%%%%%%%%%%%%%%%%%%%%%%%
\title{\textbf{The Minimum Analytic Finite Graph:\\
Truth, Individuation, and the Contact Floor}}

\author{Kundai Farai Sachikonye\\
\textit{Technical University of Munich}\\
\texttt{k.sachikonye@tum.de}}

\date{\today}
%%%%%%%%%%%%%%%%%%%%%%%%%%%%%%%%%%%%%%%%%%%%%%%%%%%%%%%%%%%%%%%%%%%%%%%%%%%%%

\begin{document}

\maketitle

%%%%%%%%%%%%%%%%%%%%%%%%%%%%%%%%%%%%%%%%%%%%%%%%%%%%%%%%%%%%%%%%%%%%%%%%%%%%%
\begin{abstract}
%%%%%%%%%%%%%%%%%%%%%%%%%%%%%%%%%%%%%%%%%%%%%%%%%%%%%%%%%%%%%%%%%%%%%%%%%%%%%

We develop a theory of individuation and truth grounded entirely in finite
weighted graphs, without appeal to set-theoretic identity, metric spaces, or
propositional logic. The central objects are: a finite weighted graph
$G = (\Vcal, \Ecal, w)$ of nodes and weighted contact edges; a distinguished
complete subgraph $\med \subseteq G$ called the \emph{medium}, which is
simultaneously in contact with every node; and the \emph{contact floor}
$\bmin > 0$, the minimum edge weight below which two nodes are
indistinguishable in $G$.

We prove three results. First, the \emph{Individuation Theorem}: a node $v
\in \Vcal$ is individuated --- distinct from the medium --- if and only if
every cut separating $v$ from the medium has weight $\geq \bmin$. Individuation
is a structural property of the graph, not a property of the node in isolation.
Second, the \emph{Truth Cell Theorem}: the truth of a node $v$ is not a
point in $\Vcal$ but the minimum cut set $C^*(v)$ of weight exactly $\bmin$
separating $v$ from the medium --- a cell, not a point. Third, the
\emph{Reshuffling Theorem}: any measurement of $v$ is a graph automorphism
(edge relabelling) that preserves the total edge weight (conservation of the
medium) and produces a new arrangement of contact edges. Measurement does not
reveal the node; it produces a new graph in which the node is individuated
differently against a new neighbour configuration. The sequence of measurements
converges not to the node's identity but to the contact floor $\bmin$.

We apply the framework to three cases: (i) population counting, where the
count of a finite population is shown to be a real-valued interval
$[n, n+1)$ at any instant of observation, not an integer; (ii) analytical
chemistry, where each instrument generation is a graph expansion adding new
neighbour nodes to the medium, increasing individuation depth without
approaching identity; (iii) the epistemology of classification, where the
sufficient truth of any classified object is the structured set of
graph-theoretic unknowables constituting its contact floor --- the minimum
weight of its boundary with the medium.

The paper contains no appeal to truth as correspondence, identity as
self-predication, or measurement as revelation. All claims are expressed
as properties of finite weighted graphs.

\end{abstract}

\vspace{1em}
\noindent\textbf{Keywords}: finite weighted graph; contact floor; individuation;
minimum cut; truth cell; graph reshuffling; medium; population counting;
analytical chemistry; epistemology of measurement; graph automorphism;
conservation

\tableofcontents
\newpage

%%%%%%%%%%%%%%%%%%%%%%%%%%%%%%%%%%%%%%%%%%%%%%%%%%%%%%%%%%%%%%%%%%%%%%%%%%%%%
\section{Introduction}
\label{sec:introduction}
%%%%%%%%%%%%%%%%%%%%%%%%%%%%%%%%%%%%%%%%%%%%%%%%%%%%%%%%%%%%%%%%%%%%%%%%%%%%%

\subsection{The Problem with Point-Truth}

Classical epistemology treats truth as a relation between a proposition and
a fact: a claim is true if it corresponds to a state of affairs. This
formulation inherits a geometric assumption: that facts are \emph{points} ---
determinate, located, self-identical. The claim ``the number of cats in the
world is $n$'' is true if and only if there is a unique integer $n$ that is
the count at the relevant moment.

This paper argues that point-truth is the wrong geometry. Not because the world
is vague or indeterminate in a philosophical sense, but because of two structural
facts that hold in any finite observable system:

\begin{enumerate}[label=(\roman*)]
  \item \textbf{Matter is conserved.} The total weight of the system does not
  change. Any rearrangement is a reshuffling of existing material, not a
  creation of new information.

  \item \textbf{Only a finite region of the system is observable at any
  instant.} Every observation is a local subgraph of a larger graph. The
  complement of the observed subgraph is always non-empty and always
  contributing to the individuation of what is observed.
\end{enumerate}

From these two facts, it follows that the truth of any node in a finite
weighted graph is not a point but a \emph{cell} --- the minimum cut set
separating the node from the rest of the graph. The floor of this cell,
the minimum achievable cut weight $\bmin > 0$, is the irreducible boundary
cost of individuation. It is not a limitation of knowledge. It is the
structural definition of what it means to be a distinct node.

\subsection{Why Graphs}

We use finite weighted graphs because they are the minimal honest language
for the claims we wish to make. A graph has:

\begin{itemize}
  \item Nodes: the things that are individuated
  \item Edges: the contact relations between things
  \item Weights: the cost of maintaining each contact relation
\end{itemize}

No metric. No measure. No probability. No propositional logic. The only
operations are: add/remove edges, relabel edges, cut the graph, count
the cut weight. Every claim in this paper is a claim about one of these
operations on a finite weighted graph.

This is not a restriction. It is a commitment to honesty about what is
actually being asserted. When we say ``the truth of $v$ is $C^*(v)$,'' we
mean: the minimum cut set separating $v$ from the medium in $G$. Nothing
more. No correspondence. No ontology beyond the graph.

\subsection{Contributions}

\begin{enumerate}[label=(\roman*)]
  \item The \emph{Medium Graph} $\med$: a distinguished complete subgraph
  simultaneously in contact with every node, playing the role of the
  invariant in which all individuation occurs
  (Section~\ref{sec:setting}).

  \item The \emph{Contact Floor} $\bmin > 0$: the minimum edge weight
  across all cuts in $G$, proved strictly positive for any finite graph
  with positive edge weights (Lemma~\ref{lem:floor_positive})
  (Section~\ref{sec:setting}).

  \item The \emph{Individuation Theorem} (Theorem~\ref{thm:individuation}):
  $v$ is individuated if and only if every cut separating $v$ from $\med$
  has weight $\geq \bmin$ (Section~\ref{sec:individuation}).

  \item The \emph{Truth Cell Theorem} (Theorem~\ref{thm:truth_cell}): the
  truth of $v$ is the minimum cut set $C^*(v)$ of weight $\bmin$ --- a
  cell, not a point (Section~\ref{sec:truth}).

  \item The \emph{Reshuffling Theorem} (Theorem~\ref{thm:reshuffling}): a
  measurement is a weight-preserving graph automorphism; it produces a new
  individuating arrangement without revealing identity; the sequence of
  measurements converges to $\bmin$, not to any node
  (Section~\ref{sec:reshuffling}).

  \item Three applications: population counting, analytical chemistry, and
  classification epistemology (Section~\ref{sec:applications}).
\end{enumerate}

%%%%%%%%%%%%%%%%%%%%%%%%%%%%%%%%%%%%%%%%%%%%%%%%%%%%%%%%%%%%%%%%%%%%%%%%%%%%%
\section{Mathematical Setting}
\label{sec:setting}
%%%%%%%%%%%%%%%%%%%%%%%%%%%%%%%%%%%%%%%%%%%%%%%%%%%%%%%%%%%%%%%%%%%%%%%%%%%%%

\subsection{The Finite Weighted Graph}

\begin{definition}[Finite Weighted Graph]
\label{def:fwg}
A \emph{finite weighted graph} is a triple $G = (\Vcal, \Ecal, w)$ where:
\begin{itemize}
  \item $\Vcal$ is a finite non-empty set of \emph{nodes};
  \item $\Ecal \subseteq \Vcal \times \Vcal$ is a set of undirected
  \emph{edges} (with $(u,v) = (v,u)$, no self-loops);
  \item $w: \Ecal \to \reals_{> 0}$ assigns a strictly positive \emph{weight}
  to each edge.
\end{itemize}
The \emph{total weight} of $G$ is $W(G) = \sum_{e \in \Ecal} w(e)$.
\end{definition}

\begin{definition}[Cut and Cut Weight]
\label{def:cut}
For a subset $S \subsetneq \Vcal$ with $S \neq \emptyset$, the \emph{cut}
$\cut(S, \bar{S})$ is the set of edges with one endpoint in $S$ and one in
$\bar{S} = \Vcal \setminus S$:
\[
  \cut(S, \bar{S}) = \{(u,v) \in \Ecal : u \in S,\, v \in \bar{S}\}.
\]
The \emph{cut weight} is $w(\cut(S,\bar{S})) = \sum_{e \in \cut(S,\bar{S})} w(e)$.
\end{definition}

\begin{definition}[Contact Floor]
\label{def:floor}
The \emph{contact floor} of $G$ is:
\[
  \bmin(G) = \min_{e \in \Ecal} w(e).
\]
It is the weight of the cheapest edge in $G$ --- the minimum cost of any
contact relation.
\end{definition}

\begin{lemma}[Positivity of the Contact Floor]
\label{lem:floor_positive}
For any finite weighted graph $G$ with $|\Ecal| \geq 1$, we have $\bmin(G) > 0$.
\end{lemma}

\begin{proof}
By Definition~\ref{def:fwg}, all edge weights are strictly positive:
$w(e) > 0$ for all $e \in \Ecal$. The minimum of finitely many strictly
positive numbers is strictly positive. Hence $\bmin = \min_{e \in \Ecal}
w(e) > 0$. \qed
\end{proof}

\begin{remark}
Lemma~\ref{lem:floor_positive} is elementary, but its content is non-trivial:
it asserts that in any finite graph with positive contact costs, there is no
zero-cost boundary. Any two distinct nodes that are connected are connected
at cost $\geq \bmin > 0$. This is the graph-theoretic statement that
individuation is never free.
\end{remark}

\subsection{The Medium}

\begin{definition}[Medium]
\label{def:medium}
The \emph{medium} of $G$ is a distinguished node $\med \in \Vcal$ (or
distinguished subset $\med \subseteq \Vcal$) that is adjacent to every
other node in $\Vcal$. Formally, for all $v \in \Vcal \setminus \{\med\}$,
$(v, \med) \in \Ecal$. The medium is the node that is simultaneously in
contact with everything.
\end{definition}

\begin{remark}[The Medium as Invariant]
\label{rem:medium_invariant}
The medium does not change when the graph is reshuffled. Its total weight
$\sum_{v \neq \med} w(v, \med)$ is conserved under any weight-preserving
automorphism (Definition~\ref{def:reshuffling} below). It is the invariant
of the system --- the thing that does not change while everything else
rearranges. In physical terms: the medium is the water in which the marbles
sit. The water is in contact with every marble. When the marbles are
reshuffled, the total volume of water is unchanged.
\end{remark}

\begin{definition}[Individuation Graph]
\label{def:indgraph}
The \emph{individuation graph} of $v \in \Vcal \setminus \{\med\}$ is the
subgraph $G_v \subseteq G$ induced by $v$, $\med$, and all nodes $u \neq v$
adjacent to both $v$ and $\med$:
\[
  \Vcal(G_v) = \{v\} \cup \{\med\} \cup \Nbd(v) \cap \Nbd(\med)
\]
where $\Nbd(u)$ denotes the neighbourhood of $u$ in $G$. The individuation
graph encodes all nodes that simultaneously individuate $v$ (by being
neighbours of $v$) and are held in the medium (by being neighbours of $\med$).
\end{definition}

\begin{remark}[Every Neighbour Contributes]
\label{rem:collective}
The individuation of $v$ is not a property of $v$ alone. It is a property
of $G_v$ --- the subgraph containing $v$, the medium, and all their shared
neighbours. Every node $u \in \Nbd(v) \cap \Nbd(\med)$ contributes an edge
$(v,u)$ to the boundary of $v$ and an edge $(u, \med)$ to the medium. The
boundary of $v$ is collectively maintained by all its neighbours. This is
the graph-theoretic content of: \emph{every other marble contributes to the
individuation of the single marble.}
\end{remark}

%%%%%%%%%%%%%%%%%%%%%%%%%%%%%%%%%%%%%%%%%%%%%%%%%%%%%%%%%%%%%%%%%%%%%%%%%%%%%
\section{Individuation}
\label{sec:individuation}
%%%%%%%%%%%%%%%%%%%%%%%%%%%%%%%%%%%%%%%%%%%%%%%%%%%%%%%%%%%%%%%%%%%%%%%%%%%%%

\subsection{What It Means to Be Distinct}

In a finite weighted graph, two nodes are \emph{the same} --- indistinguishable
--- if and only if every path between them has zero cut weight. Since all
edge weights are positive (Lemma~\ref{lem:floor_positive}), this can never
happen for adjacent nodes. The relevant question is: how much does it cost
to separate $v$ from the medium?

\begin{definition}[Separation Cost]
\label{def:sepcost}
The \emph{separation cost} of node $v$ from the medium $\med$ is:
\[
  \sigma(v) = \min_{S \ni v,\, \med \notin S} w(\cut(S, \bar{S}))
\]
the weight of the minimum cut separating $v$ from $\med$.
\end{definition}

\begin{definition}[Individuation]
\label{def:individuation}
Node $v$ is \emph{individuated} in $G$ if $\sigma(v) \geq \bmin(G)$.
Node $v$ is \emph{dissolved} if $\sigma(v) < \bmin(G)$ --- its boundary
with the medium costs less than the cheapest edge in the graph, which is
impossible for connected nodes (by Lemma~\ref{lem:floor_positive}). Therefore,
in any connected finite weighted graph, every node is individuated.
\end{definition}

\begin{theorem}[Individuation Theorem]
\label{thm:individuation}
Let $G = (\Vcal, \Ecal, w)$ be a finite connected weighted graph with medium
$\med$. Then every node $v \in \Vcal \setminus \{\med\}$ is individuated, and
its individuation depends on the full neighbourhood structure $G_v$:
\[
  \sigma(v) = \min_{S \ni v,\, \med \notin S} \sum_{u \in S,\, u' \in \bar{S}}
  w(u, u') \;\geq\; \bmin(G) > 0.
\]
Specifically, $\sigma(v)$ increases when new neighbours are added to $\Nbd(v)
\cap \Nbd(\med)$ with positive edge weights, and $\sigma(v) \geq \bmin$ always.
\end{theorem}

\begin{proof}
Connectivity of $G$ implies that at least one path connects $v$ to $\med$, so
at least one edge must be cut in any separation. The cut weight is a sum of
positive edge weights, bounded below by the minimum edge weight $\bmin > 0$
(Lemma~\ref{lem:floor_positive}). Hence $\sigma(v) \geq \bmin > 0$, so $v$
is individuated.

For the dependence on $G_v$: any node $u \in \Nbd(v) \cap \Nbd(\med)$ with
$u \notin S$ in the minimum cut contributes $w(v,u)$ to the cut weight. Adding
a new neighbour $u'$ to $\Nbd(v) \cap \Nbd(\med)$ adds either a new term
$w(v,u')$ to the cut (if $u' \notin S$) or extends $S$ without reducing the
cut weight below what any remaining cut edge contributes. In either case,
$\sigma(v)$ does not decrease. The minimum cut is collectively determined by
all neighbours. \qed
\end{proof}

\begin{corollary}[Individuation Is Collective]
\label{cor:collective}
The individuation of $v$ cannot be computed from $v$ alone. It requires
knowledge of the full neighbourhood $\Nbd(v) \cap \Nbd(\med)$ and the
weights of all edges in $G_v$. Removing the medium $\med$ from $G$ makes
$\sigma(v)$ undefined: there is no reference against which $v$ is individuated.
\end{corollary}

\begin{proof}
$\sigma(v)$ is defined as the minimum cut separating $v$ from $\med$. If
$\med$ is removed, the separation cost has no target; the concept of
individuation requires the medium as its reference. \qed
\end{proof}

%%%%%%%%%%%%%%%%%%%%%%%%%%%%%%%%%%%%%%%%%%%%%%%%%%%%%%%%%%%%%%%%%%%%%%%%%%%%%
\section{Truth as a Cell}
\label{sec:truth}
%%%%%%%%%%%%%%%%%%%%%%%%%%%%%%%%%%%%%%%%%%%%%%%%%%%%%%%%%%%%%%%%%%%%%%%%%%%%%

\subsection{The Minimum Cut Set}

\begin{definition}[Minimum Cut Set of $v$]
\label{def:mincutset}
The \emph{minimum cut set} of $v$ in $G$ is:
\[
  C^*(v) = \cut(S^*(v),\, \overline{S^*(v)})
\]
where $S^*(v) = \arg\min_{S \ni v,\, \med \notin S} w(\cut(S, \bar{S}))$
is the minimum cut partition placing $v$ on one side and $\med$ on the other.
The minimum cut set $C^*(v)$ is a set of edges, not a node. It is the
boundary of $v$ at minimum cost.
\end{definition}

\begin{remark}
$C^*(v)$ may not be unique: there may be multiple minimum cut partitions
with the same cut weight $\sigma(v)$. The \emph{truth cell} of $v$ is any
one of these minimum cut sets, or equivalently, the equivalence class of all
minimum cut partitions of $v$ from $\med$.
\end{remark}

\begin{theorem}[Truth Cell Theorem]
\label{thm:truth_cell}
The truth of node $v$ in $G$ --- the most complete structural description of
$v$ that is available within $G$ --- is not a point in $\Vcal$ but the
minimum cut set $C^*(v)$: a set of edges of total weight $\sigma(v) \geq
\bmin > 0$.

Specifically:
\begin{enumerate}[label=(\roman*)]
  \item $C^*(v)$ has positive weight: $w(C^*(v)) = \sigma(v) \geq \bmin > 0$.
  Truth is never a point (zero-weight cut).

  \item $C^*(v)$ is a boundary, not a node: it consists of edges
  connecting $S^*(v)$ to $\overline{S^*(v)}$, none of which is $v$ itself.

  \item $C^*(v)$ is collectively determined: each edge in $C^*(v)$ involves
  a neighbour $u \in \Nbd(v)$, and $u \in \Nbd(\med)$. The truth of $v$
  is held open by its neighbours, not by $v$.

  \item Increasing the graph (adding new nodes to $\Nbd(v) \cap \Nbd(\med)$)
  does not change the nature of $C^*(v)$: it remains a minimum cut set.
  Adding more neighbours makes $v$ more individuated (Theorem~\ref{thm:individuation})
  but does not collapse $C^*(v)$ to a point.
\end{enumerate}
\end{theorem}

\begin{proof}
(i) Follows directly from Lemma~\ref{lem:floor_positive} and
Definition~\ref{def:sepcost}: $w(C^*(v)) = \sigma(v) \geq \bmin > 0$.

(ii) $C^*(v)$ is a set of edges by Definition~\ref{def:mincutset}. No node
is an edge; $v \in \Vcal$ is not an element of $C^*(v) \subseteq \Ecal$.

(iii) Each edge $(u, u') \in C^*(v)$ has $u \in S^*(v)$ and $u' \in
\overline{S^*(v)}$. If $u = v$, then the edge is $(v, u')$ for some
$u' \in \Nbd(v)$; if $u \neq v$ then $u \in \Nbd(v) \cap S^*(v)$ and
$u' \in \overline{S^*(v)} \ni \med$, so $u' \in \Nbd(\med)$. In either
case, the edge involves a neighbour of $v$ in relation to the medium.

(iv) Adding node $u'$ to $\Nbd(v) \cap \Nbd(\med)$ either adds a new
edge to the minimum cut or increases $\sigma(v)$. In either case the minimum
cut remains a set of edges of positive total weight. It does not collapse. \qed
\end{proof}

\subsection{Why Truth Cannot Be a Point}

\begin{corollary}[No Point Truth]
\label{cor:no_point}
In any finite connected weighted graph with $|\Vcal| \geq 2$, the truth of
any node $v$ cannot be a point. A point-truth would require $w(C^*(v)) = 0$,
which contradicts $\bmin > 0$.
\end{corollary}

\begin{proof}
A point-truth for $v$ would mean $v$ is completely described by itself alone,
with no residual boundary cost: $\sigma(v) = 0$. But $\sigma(v) \geq \bmin > 0$
by Lemma~\ref{lem:floor_positive} applied to the minimum cut. Hence point-truth
is structurally impossible in any finite weighted graph. \qed
\end{proof}

\begin{corollary}[Truth Is Irreducibly Relational]
\label{cor:relational}
The truth of $v$ is not a property of $v$. It is a property of the minimum
cut $C^*(v)$, which is a set of edges between $v$'s neighbourhood and the
rest of $G$. Remove all edges incident to $v$ and the truth of $v$ is
undefined: $\sigma(v)$ has no cut to minimise over.
\end{corollary}

\subsection{The Unknowable Residue}

\begin{definition}[Unknowable Residue]
\label{def:unknowable}
The \emph{unknowable residue} of $v$ is the set of nodes in $S^*(v)
\setminus \{v\}$ --- the nodes that are on $v$'s side of the minimum cut
but are not $v$ itself. These nodes are indistinguishable from $v$ at the
resolution of $C^*(v)$: their separation from $v$ costs less than a minimum
cut from $\med$.
\end{definition}

\begin{proposition}[The Residue Is Structural]
\label{prop:residue}
The unknowable residue of $v$ is not a gap in knowledge. It is a structural
feature of $G$: it consists of exactly those nodes $u \neq v$ such that the
minimum cut separating $\{u, v\}$ jointly from $\med$ has the same weight as
$C^*(v)$. These nodes co-individuate with $v$ at the contact floor.
\end{proposition}

\begin{proof}
A node $u \in S^*(v) \setminus \{v\}$ is in the unknowable residue if and only
if moving $u$ to $\bar{S}$ does not increase the cut weight below $\sigma(v)$.
This means $w(u, \bar{S}) \leq 0$ at the margin --- impossible since all
weights are positive. So $u$ is in $S^*(v)$ because moving it out would
\emph{increase} the cut weight, meaning $u$'s contribution to the cut from
$v$'s side is net-negative. Equivalently: $u$ and $v$ are so tightly connected
(relative to their connections to $\med$) that they share a minimum cut
boundary. They are jointly individuated. \qed
\end{proof}

%%%%%%%%%%%%%%%%%%%%%%%%%%%%%%%%%%%%%%%%%%%%%%%%%%%%%%%%%%%%%%%%%%%%%%%%%%%%%
\section{Measurement as Graph Reshuffling}
\label{sec:reshuffling}
%%%%%%%%%%%%%%%%%%%%%%%%%%%%%%%%%%%%%%%%%%%%%%%%%%%%%%%%%%%%%%%%%%%%%%%%%%%%%

\subsection{Definition of Reshuffling}

\begin{definition}[Graph Reshuffling]
\label{def:reshuffling}
A \emph{reshuffling} of $G = (\Vcal, \Ecal, w)$ is a bijection
$\phi: \Ecal \to \Ecal'$ where $\Ecal'$ is a new edge set on $\Vcal$,
such that:
\begin{enumerate}[label=(\roman*)]
  \item \emph{Weight conservation}: $\sum_{e \in \Ecal} w(e) = \sum_{e'
  \in \Ecal'} w'(e')$. The total weight is preserved.
  \item \emph{Medium preservation}: the medium $\med$ remains adjacent to
  all nodes in $\Vcal \setminus \{\med\}$ in $G' = (\Vcal, \Ecal', w')$.
  \item \emph{Node conservation}: $\Vcal$ is unchanged.
\end{enumerate}
The reshuffled graph is $G' = (\Vcal, \Ecal', w')$.
\end{definition}

\begin{remark}[What Reshuffling Is and Is Not]
A reshuffling changes which nodes are adjacent and at what weight, while
preserving the total weight and the node set. It does not add or remove
nodes. It does not change the total weight of the medium's edges (since
the medium remains adjacent to all nodes). It produces a new individuation
--- a different assignment of neighbourhoods --- without creating new matter.

This is the graph-theoretic statement of: \emph{matter is conserved, and
reshuffling produces new arrangements of surfaces and volumes, not new matter.}
\end{remark}

\begin{theorem}[Reshuffling Theorem]
\label{thm:reshuffling}
Let $\phi$ be a reshuffling of $G$, producing $G'$. Then:
\begin{enumerate}[label=(\roman*)]
  \item The contact floor is preserved or increased: $\bmin(G') \geq \bmin(G)$.
  \item The truth cell of $v$ in $G'$ is $C^*_{G'}(v)$, which may differ from
  $C^*_G(v)$: the reshuffling produces a new individuation of $v$.
  \item The separation cost $\sigma_{G'}(v)$ satisfies $\sigma_{G'}(v) \geq
  \bmin(G) > 0$: $v$ remains individuated after reshuffling.
  \item No finite sequence of reshufficings $\phi_1, \phi_2, \ldots, \phi_k$
  produces a graph in which $\sigma(v) = 0$: the floor is never reached by
  reshuffling.
\end{enumerate}
\end{theorem}

\begin{proof}
(i) The minimum edge weight in $G'$ is $\min_{e' \in \Ecal'} w'(e')$. Since
total weight is conserved and the number of edges may change, the minimum
weight per edge can only increase (if edges are merged) or stay the same.
Specifically, reshuffling cannot create an edge of weight less than
$\bmin(G)$ without removing a heavier edge, which would require weight
creation --- forbidden by conservation. Hence $\bmin(G') \geq \bmin(G)$.

(ii) In $G'$, the neighbourhood $\Nbd_{G'}(v)$ differs from $\Nbd_G(v)$
(different edges). The minimum cut $C^*_{G'}(v)$ is recomputed in the new
graph and generally differs from $C^*_G(v)$.

(iii) By the same argument as Theorem~\ref{thm:individuation} applied to
$G'$: connectivity of $G'$ (which follows from medium preservation) and
positivity of edge weights give $\sigma_{G'}(v) \geq \bmin(G') \geq
\bmin(G) > 0$.

(iv) By (iii), $\sigma(v) \geq \bmin(G) > 0$ after every reshuffling. No
finite composition of reshufficings reduces $\sigma(v)$ to zero. The sequence
$(\sigma_k(v))_{k \geq 0}$ is bounded below by $\bmin(G) > 0$. \qed
\end{proof}

\subsection{The Convergence of Measurement}

\begin{corollary}[Measurement Converges to the Floor, Not to a Node]
\label{cor:convergence}
Let $(\phi_k)_{k \geq 1}$ be a sequence of reshufficings of $G$ such that
$\bmin(G_k) \to L$ as $k \to \infty$. Then:
\begin{enumerate}[label=(\roman*)]
  \item $L \geq \bmin(G_0) > 0$: the limit is positive.
  \item $\sigma_k(v) \to L > 0$: the separation cost converges to the
  floor, not to zero.
  \item The truth cell $C^*_{G_k}(v)$ converges (in the sense of its total
  weight) to a cut of weight $L > 0$: a cell, not a point.
\end{enumerate}
In particular, no sequence of measurements can produce a truth cell of zero
weight. The node $v$ is never fully revealed by measurement.
\end{corollary}

\begin{proof}
By Theorem~\ref{thm:reshuffling}(i), $\bmin(G_k)$ is non-decreasing and
bounded. It therefore converges to some $L \geq \bmin(G_0) > 0$. By (iii),
$\sigma_k(v) \geq \bmin(G_k) \to L > 0$. The weight of the truth cell is
$w(C^*_{G_k}(v)) = \sigma_k(v) \to L > 0$. \qed
\end{proof}

\subsection{What Measurement Produces}

\begin{proposition}[Measurement as NOT-Sequence]
\label{prop:not_sequence}
Each reshuffling $\phi_k$ produces a new graph $G_k$ in which $v$ has a new
neighbourhood $\Nbd_{G_k}(v)$. The new neighbours are the nodes that
participate in individuating $v$ in the new arrangement. The information
content of the sequence $(\phi_1, \ldots, \phi_k)$ is not ``what $v$ is''
but ``what $v$ is not, relative to each neighbourhood configuration.''

Formally, the sequence produces a sequence of minimum cut sets:
\[
  C^*_{G_0}(v),\; C^*_{G_1}(v),\; \ldots,\; C^*_{G_k}(v)
\]
each of which is a set of NOT-edges: edges between $v$'s current partition
and its complement. The intersection $\bigcap_{k} \mathrm{supp}(C^*_{G_k}(v))$
--- the edges that appear in every minimum cut set --- is the irreducible
boundary of $v$: the contact floor $\bmin$.
\end{proposition}

\begin{proof}
Each $C^*_{G_k}(v)$ is a minimum cut set in $G_k$. By
Corollary~\ref{cor:convergence}, its weight converges to $L \geq \bmin(G_0)$.
The edges in $C^*_{G_k}(v)$ are those that must be cut in any minimum
separation of $v$ from $\med$ in $G_k$. These are precisely the NOT-edges:
edges that go between $v$'s region and the rest. As $k \to \infty$, the
residual common structure in all cuts is the floor $\bmin$. \qed
\end{proof}

%%%%%%%%%%%%%%%%%%%%%%%%%%%%%%%%%%%%%%%%%%%%%%%%%%%%%%%%%%%%%%%%%%%%%%%%%%%%%
\section{Applications}
\label{sec:applications}
%%%%%%%%%%%%%%%%%%%%%%%%%%%%%%%%%%%%%%%%%%%%%%%%%%%%%%%%%%%%%%%%%%%%%%%%%%%%%

\subsection{Population Counting}

\begin{example}[The Cat Counter]
\label{ex:cats}
Suppose we wish to determine the number of cats in the world. We are given:
a strict classification $\mathcal{C}$ of what constitutes a cat (a finite
set of defining properties); a finite set of laser counters that register
transitions (births and deaths) in real time.

Model this as a finite weighted graph $G$ where:
\begin{itemize}
  \item Each cat is a node $v_i \in \Vcal$;
  \item The medium $\med$ is the world-state (the invariant in which all
  cats exist);
  \item Edge weights $w(v_i, \med) = 1$ for each living cat (the cost of
  individuation from the world), $w(v_i, \med) = 0$ for a dead cat --- but
  $w = 0$ is forbidden by Lemma~\ref{lem:floor_positive}, meaning the node
  is removed from $\Vcal$ (death = removal, not zero-weight edge).
  \item Each birth is a node addition; each death is a node removal.
\end{itemize}

The count $|\Vcal \setminus \{\med\}|$ is the number of cat-nodes. But:
the laser counter registers transitions. At any instant $t$, there is a
non-zero probability that a transition is in progress (a cat is being born
or dying). The count at $t$ is therefore not a single integer but the
interval $[n(t), n(t)+1)$ where $n(t)$ is the last confirmed count.

Formally: the total weight $W(G_t) = \sum_i w(v_i, \med)$ is conserved
(matter conservation). Node additions/removals change $|\Vcal|$ but not
$W(G)$ instantaneously --- the weight must be redistributed from the
removed node's edges to the remaining graph. During a transition, the
graph $G_t$ is in a reshuffling state: the count is simultaneously $n$ and
$n+1$ in the sense that both partitions of the graph ($n$ cat-nodes and
$n+1$ cat-nodes) have the same total weight.
\end{example}

\begin{theorem}[Count as Interval]
\label{thm:count_interval}
For any finite population modelled as a finite weighted graph with continuous
transitions, the count at any instant $t$ is an element of the interval
$[n(t), n(t)+1)$, not a single integer. The binary partition
$\{n, n+1\}$ is the finest partition of the count, and the population is
simultaneously even and odd during any transition.
\end{theorem}

\begin{proof}
At a birth event, a new node $v_{\mathrm{new}}$ is added to $\Vcal$ and
connected to $\med$ with weight $\bmin$. The transition is not instantaneous:
the edge $(v_{\mathrm{new}}, \med)$ acquires weight continuously from $0$ to
$\bmin$. During the interval when $w(v_{\mathrm{new}}, \med) \in (0, \bmin)$,
the node exists (positive weight, hence individuated at a supraminimal cut)
but has not yet reached the floor cost of full individuation. The count is
therefore not $n$ (the new node is present) and not $n+1$ (the floor has not
been reached). It is in $[n, n+1)$.

The binary partition $\{$even, odd$\}$ applied to this interval gives both
$n$ (if even) and $n+1$ (if odd) as valid descriptions simultaneously.
Neither alone is sufficient. \qed
\end{proof}

\begin{remark}[The Best Available Truth]
\label{rem:count_truth}
The even/odd partition is the coarsest binary partition of the count. It is
also the most honest: it acknowledges that the count is in $[n, n+1)$ by
assigning it to both halves. The minimum sufficient truth of the population
count is the pair $\{n, n+1\}$ --- not a single integer, but the smallest
interval consistent with the graph structure at time $t$.
\end{remark}

\subsection{Analytical Chemistry: Instrument Generations as Graph Expansion}

\begin{example}[The Mass Spectrometer as Graph Expansion]
\label{ex:ms}
An analytical instrument that separates a mixture into components models as
a finite weighted graph where:
\begin{itemize}
  \item Each distinguishable component species is a node $v_i$;
  \item The medium $\med$ is the solvent/matrix/background --- the invariant
  simultaneously in contact with all species;
  \item Edge weights $w(v_i, v_j)$ encode the cost of distinguishing species
  $i$ from species $j$ (the resolving effort required);
  \item The contact floor $\bmin$ is the minimum resolving cost across all
  pairs --- the finest distinction the instrument can make.
\end{itemize}

A measurement is a reshuffling $\phi$: the instrument rearranges which species
are compared to which, producing a new neighbourhood for each node.
\end{example}

\begin{theorem}[Instrument Generation as Neighbour Expansion]
\label{thm:instrument_generation}
Each new generation of analytical instrument corresponds to a graph expansion:
the addition of new nodes to $\Nbd(v) \cap \Nbd(\med)$ for every node $v$.
The effect is:
\begin{enumerate}[label=(\roman*)]
  \item The individuation of each node increases (Theorem~\ref{thm:individuation}):
  $\sigma_{G'}(v) \geq \sigma_G(v)$.
  \item The truth cell $C^*(v)$ changes: new edges appear in the minimum
  cut, corresponding to new NOT-relations established by the new neighbours.
  \item The contact floor $\bmin$ does not decrease: the minimum resolving
  cost is preserved or improved.
  \item No finite number of instrument generations collapses $C^*(v)$ to a
  point: the truth cell remains a cell at every generation.
\end{enumerate}
\end{theorem}

\begin{proof}
(i) By Theorem~\ref{thm:individuation}, adding neighbours to $\Nbd(v) \cap
\Nbd(\med)$ does not decrease $\sigma(v)$.

(ii) The new neighbours $u' \in \Nbd_{G'}(v) \setminus \Nbd_G(v)$ contribute
new edges to the cut $C^*_{G'}(v)$: the NOT-relation $v \neq u'$ is now an
edge in the boundary.

(iii) By Theorem~\ref{thm:reshuffling}(i), $\bmin(G') \geq \bmin(G)$.

(iv) By Corollary~\ref{cor:convergence}, the truth cell weight converges to
$L \geq \bmin > 0$. It never reaches zero. \qed
\end{proof}

\begin{corollary}[What Each Generation Produces]
\label{cor:generation}
Each instrument generation does not produce more truth about $v$. It produces
a larger NOT-sequence: a longer list of nodes from which $v$ has been
distinguished. The truth cell $C^*(v)$ grows in the number of its edges
(more NOT-relations) but not in its weight towards zero. Resolving power is
neighbour count, not proximity to identity.
\end{corollary}

\begin{remark}[The Sample Preparation Regress]
\label{rem:regress}
The analytical result depends on the sample inserted into the instrument.
Confirming the sample's composition requires a prior measurement, which
requires a prior confirmed sample, and so on. In graph terms: the node
$v_{\mathrm{sample}}$ representing the inserted sample is itself individuated
only within the graph of all prior preparation steps. Each preparation step
is a reshuffling of the upstream graph. The contact floor of the entire
preparation chain is the accumulated floor across all upstream graphs:
$\bmin^{\mathrm{chain}} = \min_k \bmin(G_k) > 0$. The floor of the
result is bounded below by the floor of the coarsest upstream graph,
not by the instrument alone.
\end{remark}

\subsection{The Epistemology of Classification}

\begin{example}[Classification as Graph Partition]
\label{ex:classification}
A classification system assigns nodes to categories. Each category $C_j$
is a subset of $\Vcal$. The classification is a partition $\Pcal =
\{C_1, \ldots, C_k\}$ of $\Vcal \setminus \{\med\}$. The contact floor
$\bmin(\Pcal)$ of the classification is the minimum cut weight between any
two categories across the medium.
\end{example}

\begin{theorem}[Sufficient Truth of Classification]
\label{thm:classification_truth}
The sufficient truth of a classification $\Pcal$ of $G$ is not the partition
itself but the contact floor $\bmin(\Pcal) > 0$: the minimum cost at which
any two categories can be distinguished through the medium. The sufficient
truth is a positive number, not a partition assignment.
\end{theorem}

\begin{proof}
The partition $\Pcal$ assigns each node to a category. But the assignment
is only meaningful insofar as the categories are distinguishable --- i.e.,
insofar as the cut between any two categories has positive weight. The
minimum such weight is $\bmin(\Pcal)$. If $\bmin(\Pcal) = 0$, two categories
are indistinguishable (they share a zero-weight boundary) and the classification
provides no information. Since $\bmin(\Pcal) > 0$ for any finite weighted
graph, every classification provides some information --- but the quantity of
information is exactly $\bmin(\Pcal)$: the minimum cost of any distinction
the classification makes. This is the sufficient truth. \qed
\end{proof}

\begin{corollary}[Sufficient Truth Is the Floor, Not the Label]
\label{cor:sufficient_truth}
The sufficient truth of assigning $v$ to category $C_j$ is not the label
$C_j$. It is the minimum cut weight separating $C_j$ from all other
categories: $\bmin(\{C_j, \overline{C_j}\}) \geq \bmin(\Pcal) > 0$.
The label is a convenience. The floor is the content.
\end{corollary}

%%%%%%%%%%%%%%%%%%%%%%%%%%%%%%%%%%%%%%%%%%%%%%%%%%%%%%%%%%%%%%%%%%%%%%%%%%%%%
\section{Discussion}
\label{sec:discussion}
%%%%%%%%%%%%%%%%%%%%%%%%%%%%%%%%%%%%%%%%%%%%%%%%%%%%%%%%%%%%%%%%%%%%%%%%%%%%%

\subsection{The Medium as the Only Invariant}

Throughout the paper, the medium $\med$ is the single invariant. It does not
change under reshuffling. Its total edge weight is conserved. It is the
condition under which individuation occurs. Without the medium, no node can
be individuated --- $\sigma(v)$ is undefined without a reference.

This is the graph-theoretic statement of a physical fact: in a closed system,
the total weight (mass, energy, information) is conserved. Reshuffling
rearranges how that weight is distributed across edges, but the total is fixed.
The medium is what holds the total fixed while the distribution changes.

\subsection{Truth as Region}

The Truth Cell Theorem establishes that truth in a finite weighted graph is
always a cell --- a set of edges of positive total weight --- and never a
point. This is not a limitation of knowledge or measurement. It is a structural
consequence of $\bmin > 0$: any system in which individuation is possible
(any system with positive edge weights) has irreducibly cellular truth.

The cell is the minimum viable boundary. Below the cell, the node merges
with the medium: it is no longer distinct. The cell is not the gap between
current knowledge and eventual identity. It is the structure of identity
itself. Identity, in a finite weighted graph, is always a boundary, never
a point.

\subsection{Reshuffling and the Arrow of Knowledge}

The Reshuffling Theorem establishes that measurement does not reduce the
truth cell to a point. Each measurement produces a new arrangement of
NOT-relations --- a new set of neighbours from which the node has been
distinguished. The sequence of measurements converges to the floor $\bmin$,
not to zero.

This gives an \emph{arrow of knowledge}: the sequence of NOT-sets grows
monotonically (new NOT-relations are never removed, only added), and the
floor is the asymptote. The direction of the arrow is determined by
$\bmin > 0$: knowledge cannot exceed the floor.

This is the graph-theoretic counterpart of the physical arrow of time:
just as inflation cannot reverse below the contact floor in a dynamical
system, knowledge cannot deepen below $\bmin$ in a finite weighted graph.
The floor is the common invariant.

\subsection{Relation to Minimum Cut Theory}

The minimum cut problem in graph theory asks: what is the minimum weight
set of edges whose removal disconnects $G$? \citep{ford1956maxflow,
stoer1997simple}. The Truth Cell Theorem repurposes this: instead of
asking what removes connectivity, we ask what maintains the minimum
connection between $v$ and $\med$. The minimum cut separating $v$ from
$\med$ is the minimum cost boundary of $v$'s individuation.

The max-flow min-cut theorem \citep{ford1956maxflow} gives an algorithmic
route to computing $\sigma(v) = C^*(v)$: it equals the maximum flow from
$v$ to $\med$ in $G$, which is the maximum rate at which individuation
information can flow from $v$ to the medium. The truth cell weight is the
information bottleneck.

\subsection{Relation to Information Theory}

The contact floor $\bmin$ is related to the minimum description length of
a node in $G$: the smallest number of bits needed to specify the position
of $v$ in the minimum cut $C^*(v)$ against the medium. By Shannon's theorem
\citep{shannon1948mathematical}, this is bounded below by $\log_2(1/p_v)$
where $p_v$ is the probability of $v$'s minimum cut configuration. Since
$p_v < 1$ (there are always other configurations), $\log_2(1/p_v) > 0$:
the minimum description length is positive, consistent with $\bmin > 0$.

\subsection{What This Framework Does Not Say}

This paper makes no claim about:
\begin{itemize}
  \item The ontological status of nodes: whether they ``really exist''
  independently of the graph;
  \item The correspondence between graph structure and physical reality;
  \item Whether the medium corresponds to any particular physical substrate;
  \item The computability of the contact floor in any specific instance.
\end{itemize}

All claims are structural claims about finite weighted graphs. The
applications in Section~\ref{sec:applications} are modelling choices:
the claim is ``if you model the system as a finite weighted graph with
these properties, then the following graph-theoretic results hold.'' The
modelling choice is a scientific question; the graph-theoretic results are
mathematical facts.

%%%%%%%%%%%%%%%%%%%%%%%%%%%%%%%%%%%%%%%%%%%%%%%%%%%%%%%%%%%%%%%%%%%%%%%%%%%%%
\section{Conclusion}
\label{sec:conclusion}
%%%%%%%%%%%%%%%%%%%%%%%%%%%%%%%%%%%%%%%%%%%%%%%%%%%%%%%%%%%%%%%%%%%%%%%%%%%%%

We have developed a theory of individuation and truth grounded entirely in
finite weighted graphs. The main results are:

\begin{enumerate}
\item \textbf{The contact floor is strictly positive}
  (Lemma~\ref{lem:floor_positive}): in any finite graph with positive edge
  weights, there is no zero-cost boundary. Individuation is never free.

\item \textbf{Individuation is collective} (Theorem~\ref{thm:individuation}):
  the separation cost $\sigma(v)$ depends on the full neighbourhood of $v$
  through the medium. No node individuates itself. Every neighbour
  contributes.

\item \textbf{Truth is a cell} (Theorem~\ref{thm:truth_cell}): the truth of
  a node is its minimum cut set $C^*(v)$ --- a set of edges of positive
  total weight $\geq \bmin > 0$. Truth is never a point.

\item \textbf{Measurement is reshuffling} (Theorem~\ref{thm:reshuffling}):
  a measurement is a weight-preserving relabelling of edges. It produces
  a new individuation of $v$ --- a new set of NOT-relations --- without
  revealing $v$'s identity. Total weight is conserved. The medium is
  unchanged.

\item \textbf{Measurement converges to the floor}
  (Corollary~\ref{cor:convergence}): no sequence of measurements reduces
  the truth cell to a point. The sequence converges to the floor $\bmin > 0$.
  The node is never fully revealed.

\item \textbf{The sufficient truth of a classification is the floor}
  (Theorem~\ref{thm:classification_truth}): the sufficient truth of
  assigning $v$ to a category is not the category label but the minimum
  cut weight between that category and all others through the medium.
  The floor is the content; the label is the convenience.
\end{enumerate}

The contact floor $\bmin > 0$ is the single invariant from which all
results follow. It is not a measurement artifact. It is not a limit of
technology. It is the structural cost of any boundary in a finite graph
with positive edge weights --- the minimum price of distinction. Below the
floor, there is no distinction. Above the floor, there is always a cell,
never a point.

The medium is the other invariant: the thing simultaneously in contact with
everything, unchanged by reshuffling, conserved in total weight. It is not
a node like the others. It is the condition under which the others are nodes
at all.

Truth, in this framework, is the minimum cut between a node and the medium:
the least it can be while still being distinct. It is a region. It is positive.
It is collectively held open by every neighbour. And it is irreducibly
unknowable --- not because knowledge is limited, but because the boundary
itself is what the node is.

%%%%%%%%%%%%%%%%%%%%%%%%%%%%%%%%%%%%%%%%%%%%%%%%%%%%%%%%%%%%%%%%%%%%%%%%%%%%%
\bibliographystyle{plainnat}
\bibliography{references}
%%%%%%%%%%%%%%%%%%%%%%%%%%%%%%%%%%%%%%%%%%%%%%%%%%%%%%%%%%%%%%%%%%%%%%%%%%%%%

\end{document}
